\subsection{Sprint 1 Analysis}
\label{sec:static-sprint1}

Sprint 1 had no CI pipeline, and PHPMD and Pint are not part of the CI in any sprint.
For this reason, the tools were run by hand in a terminal.
The ``before'' run used the deployed code (tag \texttt{sprint1-deploy}).
The ``after'' run used a new branch with small fixes.
The branch is on GitHub, but it was not merged or deployed.

Each tool counts differently.
PHPStan counts errors, PHPMD counts violations, Pint counts files that need formatting, and ESLint counts errors and warnings.
PHPStan was run twice.
The first run uses the project baseline, as the CI does.
The second run ignores the baseline, so it shows the real debt.

\begin{table}[H]
  \centering
  \begin{tabularx}{\linewidth}{L{4.6cm}YY}
    \toprule
    & Before fixes & After fixes \\
    \midrule
    Date (UTC) & 2026-11-17 & 2026-11-17 \\
    Branch & \texttt{sprint1-deploy} (tag) & \texttt{fix/static-}\allowbreak\texttt{analysis-sprint1} \\
    Commit SHA & \sha{3bbef9d} & \sha{67c100a} \\
    \midrule
    PHPStan, with baseline & 112 errors & 0 errors \\
    PHPStan, without baseline & 144 errors & 64 errors \\
    PHPMD & 25 violations & 23 violations \\
    Pint & 6 files & 0 files \\
    ESLint, frontend & 7 errors, 3 warnings & 0 errors, 2 warnings \\
    ESLint, Leodeguita & 0 errors, 0 warnings & 0 errors, 0 warnings \\
    \bottomrule
  \end{tabularx}
  \caption{Sprint 1 static analysis results.}
  \label{tab:static-sprint1}
\end{table}

\evidencefigure{static-analysis/sprint1-before.png}{Sprint 1 analysis before fixes.}{fig:static-sprint1-before}

\evidencefigure{static-analysis/sprint1-after.png}{Sprint 1 analysis after fixes.}{fig:static-sprint1-after}

\subsubsection*{Issues by Type and Severity}

PHPStan and Pint have no severity levels, so every finding counts as an error.
PHPMD reports a priority from 1 (highest) to 5.
All findings below have priority 1 or 3.
ESLint separates errors from warnings.

\begin{table}[H]
  \centering
  \begin{tabularx}{\linewidth}{L{2.4cm}YL{1.6cm}L{1.4cm}L{1.4cm}}
    \toprule
    Tool & Type & Severity & Before & After \\
    \midrule
    PHPStan & Unknown property (\texttt{property.notFound}) & Error & 94 & 50 \\
    PHPStan & Unknown relation (\texttt{larastan.relationExistence}) & Error & 30 & 0 \\
    PHPStan & Invalid PHPDoc annotation & Error & 4 & 0 \\
    PHPStan & Other type problems & Error & 16 & 14 \\
    PHPMD & Unused parameter & P3 & 17 & 16 \\
    PHPMD & Undefined variable & P3 & 3 & 3 \\
    PHPMD & Boolean flag argument & P1 & 2 & 2 \\
    PHPMD & Complexity and class size & P3 & 2 & 2 \\
    PHPMD & Else expression & P1 & 1 & 0 \\
    Pint & Code style & Error & 6 & 0 \\
    ESLint & Unused variable or import & Error & 5 & 0 \\
    ESLint & Fast refresh export & Error & 2 & 0 \\
    ESLint & Missing hook dependency & Warning & 2 & 2 \\
    ESLint & Unused disable directive & Warning & 1 & 0 \\
    \bottomrule
  \end{tabularx}
  \caption{Sprint 1 findings by tool and type. PHPStan counts are without baseline.}
  \label{tab:static-sprint1-types}
\end{table}

\subsubsection*{Representative Fixes}

\begin{itemize}[nosep]
  \item \textbf{Pint} (\sha{7fbb639}): the formatter fixed 6 files. It changed only spacing, imports and comments.
  \item \textbf{ESLint} (\sha{fe7e93f}): removed unused imports and one unused disable directive.
        Allowed the ``unused key in a rest object'' pattern.
        Kept the rule that warns about exports in one file only for \texttt{WizardContext.tsx}, where the context, provider and hook belong together.
  \item \textbf{PHPStan} (\sha{531d3f1}): added return types to 55 model relations in 21 models.
        Together with the annotation fix below, this cut PHPStan errors from 144 to 64.
  \item \textbf{PHPMD} (\sha{62ccb6f}): removed one \texttt{else} branch and one unused parameter.
        Quoted four \texttt{@SuppressWarnings} annotations so PHPStan can read them.
  \item \textbf{Baseline} (\sha{67c100a}): regenerated the PHPStan baseline. It now lists only the 64 remaining errors.
\end{itemize}

No behavior changed.
After the fixes, the backend suite passed (851 tests) and the frontend suite passed (518 tests).

\subsubsection*{Unresolved Issues}

\begin{itemize}[nosep]
  \item \textbf{PHPStan, 64 errors.}
        They are recorded in the baseline, not fixed.
        40 come from two \texttt{Resource} classes that read \texttt{\$this->id} and similar fields.
        10 come from an attribute added at runtime in \texttt{ReservationsController}.
        The other 14 are small type problems.
        The with-baseline result (0 errors) hides this debt, so both numbers are reported.
  \item \textbf{PHPMD, 23 violations.}
        15 of the 16 unused parameters are required by the framework (casts, exception \texttt{render}, resource \texttt{toArray} and a validation closure), so they cannot be removed.
        The last one, \texttt{\$actingUserId} in \texttt{ReservationService::create()}, is kept to match the other methods of the service.
        3 undefined variable reports are false positives: \texttt{\$requiresReview} is passed by reference.
        4 are design issues (boolean flags, one complex method, one class with many public methods) and need a larger refactor.
  \item \textbf{ESLint, 2 warnings.}
        The \texttt{useEffect} hooks in \texttt{PreguntaInicio5} and \texttt{PreguntaInicio6} miss a dependency.
        Adding it could change when data is saved to \texttt{localStorage}, so they were left as they are.
  \item \textbf{Not covered.}
        Pint and PHPMD are not run by the CI, so they depend on manual runs.
\end{itemize}
